\begin{beginnerstatement}[Kurz erklärt: Konfiguration]

Eine Konfiguration ist eine Sammlung von Einstellungen, die festlegt, wie und
wo ein Programm ausgeführt wird. Eine Umgebungsvariable ist ein Wert, den die
Laufzeitumgebung dem Programm bereitstellt; eine lokale \texttt{.env}-Datei
kann solche Werte für die Entwicklung sammeln. \path{environment.toml}
beschreibt die erlaubten Variablen und mögliche Defaults, also verwendete
Standardwerte. \path{project-profile.toml} legt dagegen fest, welche
technischen Bausteine zum Projekt gehören. Ein Override ersetzt einen
niedriger priorisierten Wert bewusst, und der Resolver ermittelt daraus nach
einer festen Reihenfolge den tatsächlich geltenden Wert. Ein Secret ist eine
geheime Angabe wie ein Passwort und darf nicht an den sichtbaren Client
gelangen. \texttt{DATABASE\_URL} enthält die serverseitige
Datenbankverbindung und wird maskiert. \texttt{VITE\_API\_BASE\_URL} ist
dagegen die öffentliche Backend-Adresse, die das Frontend verwenden darf.

\end{beginnerstatement}
